% 431 项目统一文档主题。
% 各根入口应先声明 \documentclass{exam-zh}，再输入本文件。

% 项目统一中文字体：避免依赖某台电脑上的 STSong/STHeiti 等系统字体。
\setCJKmainfont{SourceHanSerifCN-Regular.ttf}[
    Path = fonts/,
    BoldFont = SourceHanSerifCN-Bold.ttf,
    ItalicFont = AR-PL-KaitiM-GB.ttf,
    BoldItalicFont = AR-PL-KaitiM-GB.ttf,
    BoldItalicFeatures = {FakeBold = 2},
]
\setCJKsansfont{SourceHanSansCN-Regular.ttf}[
    Path = fonts/,
    BoldFont = SourceHanSansCN-Bold.ttf,
]
\setCJKmonofont{SourceHanSansCN-Regular.ttf}[
    Path = fonts/,
]

% fontset=none 时 CTeX 不创建这些常用字体族命令，显式补齐。
\setCJKfamilyfont{zhsong}{SourceHanSerifCN-Regular.ttf}[Path = fonts/]
\setCJKfamilyfont{zhhei}{SourceHanSansCN-Bold.ttf}[Path = fonts/]
\setCJKfamilyfont{zhkai}{AR-PL-KaitiM-GB.ttf}[Path = fonts/]
\setCJKfamilyfont{zhfs}{SourceHanSerifCN-Regular.ttf}[Path = fonts/]
\providecommand{\songti}{\CJKfamily{zhsong}}
\providecommand{\heiti}{\CJKfamily{zhhei}}
\providecommand{\kaishu}{\CJKfamily{zhkai}}
\providecommand{\fangsong}{\CJKfamily{zhfs}}


\usepackage{booktabs}
\usepackage{multirow}
\usepackage{threeparttable}
\usepackage{threeparttablex}
\usepackage{tikz}
\usetikzlibrary{calc}
\usetikzlibrary{arrows.meta}
\usetikzlibrary{decorations.pathmorphing}
\usetikzlibrary{decorations.pathreplacing}
\usetikzlibrary{positioning}

\tikzset{
    mindmap_root/.style={rectangle, rounded corners, draw=black, fill=gray!10, thick, align=center, minimum height=0.8cm, font=\bfseries\small},
    mindmap_level1/.style={rectangle, rounded corners, draw=blue!60, fill=blue!5, thin, align=center, minimum height=0.6cm, font=\small},
    mindmap_level2/.style={rectangle, draw=none, fill=none, align=center, font=\footnotesize, text width=3cm},
    arrow/.style={->, >=stealth, thick, gray}
}

\usepackage{tabularx}
\usepackage{xcolor}
\usepackage{makecell}
\usepackage{extarrows}
\usepackage{longtable}
\usepackage{float}
\usepackage{enumitem}
\usepackage{perpage}
\usepackage{scalerel}
\usepackage{stackengine}
\MakePerPage{footnote}
\stackMath
\tcbuselibrary{breakable}
\tcbset{breakable}

% 适用于较长估计变量名的宽帽号。
\newcommand{\reallywidehat}[1]{%
  \savestack{\tmpbox}{\stretchto{%
    \scaleto{%
      \scalerel*[\widthof{\ensuremath{#1}}]{\kern-.6pt\bigwedge\kern-.6pt}{%
        \rule[-\textheight/2]{1ex}{\textheight}%
      }%
    }{\textheight}%
  }{0.5ex}}%
  \stackon[1pt]{#1}{\tmpbox}%
}

% 列表全局设置
\setlist{nosep}
% 列表内部保持正文节奏；整个列表块与上下内容各保留半行间距。
\setlist[itemize]{
    leftmargin=*, labelsep=0.5em,
    itemsep=0pt, parsep=0pt,
    topsep=0.5\baselineskip, partopsep=0pt,
    after={}
}
\setlist[itemize,1]{label=$\bullet$}
\setlist[itemize,2]{label=\small$\bullet$}
\setlist[itemize,3]{label=\scriptsize$\bullet$}
\setlist[enumerate]{
    leftmargin=*, labelsep=0.5em,
    itemsep=0pt, parsep=0pt,
    topsep=0.5\baselineskip, partopsep=0pt,
    after={}
}
\setlist[description]{leftmargin=\parindent, labelindent=0pt}

% 章节公共样式。需要每节另起一页的资料在入口中单独覆盖 section/break。
\IfFileExists{ctex.sty}{
    \ctexset{
        section = {
            format = \Large\bfseries\raggedright,
            beforeskip = 1.5ex plus 0.5ex minus 0.2ex,
            afterskip = 1.5ex plus 0.2ex,
        },
        subsection = {
            format = \large\bfseries\raggedright,
            beforeskip = 1.0ex plus 0.5ex minus 0.2ex,
            afterskip = 1.0ex plus 0.2ex,
        },
        subsubsection = {
            format = \normalsize\bfseries\raggedright,
            beforeskip = 0.5ex plus 0.5ex minus 0.2ex,
            afterskip = 0.5ex plus 0.2ex,
        },
        paragraph = {
            runin = false,
            number = \arabic{paragraph}.,
            format = \normalsize\bfseries\raggedright,
            indent = 2em,
            beforeskip = 0.5ex plus 0.3ex minus 0.1ex,
            afterskip = 0.3ex plus 0.2ex,
        }
    }
}{}

% 按 chapter 计数器自动生成“第 X 讲”。
% exam-zh 在 \chapter 完成时推进 chapter，因此标题使用当前值加一。
\ExplSyntaxOn
\NewDocumentCommand{\lecture}{m}
  {
    \exp_args:Nx \chapter
      {第\zhnumber{\int_eval:n {\value{chapter} + 1}}讲：\exp_not:n {#1}}
  }
\ExplSyntaxOff

% 431 资料通用的 exam-zh 设置；可选参数为居中页眉，必选参数为右侧页眉。
\newcommand{\SetupFourThreeOneDocument}[2][]{
    \examsetup{
        page={
            size = a4paper,
            show-head = true,
            head-content = {
                \fancyhf{}
                \fancyhead[C]{#1}
                \fancyhead[EL]{\leftmark}
                \fancyhead[OR]{#2}
            },
            foot-content = {\hfill 第;页 / 共;页 \hfill}
        },
        style/student-version-cleanaux = true,
        fillin/type = paren,
        question = {
            show-points = auto,
            show-answer = true,
        },
        problem = {
            points-separate-par = false,
        },
        paren = {
            text-color = blue,
            show-paren = true,
            show-answer = true,
        },
        solution = {
            show-solution = show-stay,
            label-indentation = true,
            text-color = blue,
            blank-type = hide,
        },
    }

    \ExamPrintAnswerSet{
        question/show-answer = false,
        paren/show-answer = false,
        solution/show-solution = hide,
        paren/text-color = black,
        solution/text-color = black,
    }
}

\renewcommand{\chaptermark}[1]{\markboth{#1}{}}
\renewcommand{\headrulewidth}{1pt}

\newcommand{\true}{$\checkmark$}
\newcommand{\false}{$\times$}

\usepackage{hyperref}
\usepackage{tocloft}

\renewcommand{\contentsname}{\hspace*{\fill}目录\hspace*{\fill}}
\setlength{\cftbeforetoctitleskip}{2em}
\setlength{\cftaftertoctitleskip}{2em}
\setlength{\cftchapindent}{0em}
\setlength{\cftsecindent}{2.0em}
\setlength{\cftsecnumwidth}{2em}
\setlength{\cftsubsecindent}{4.0em}
\setlength{\cftsubsecnumwidth}{2.0em}
\setlength{\cftsubsubsecindent}{6.0em}
\setlength{\cftsubsubsecnumwidth}{2.0em}

\setlength{\intextsep}{1em}

\hypersetup{
    colorlinks = true,
    linkcolor = blue,
    urlcolor = blue,
    citecolor = blue,
    pdfborder = {0 0 0},
}

\newcommand{\kref}[1]{\ref{#1} \textsuperscript{(p.\pageref{#1})}}
